% ===== PRÉAMBULE CANONIQUE — à coller VERBATIM en tête de CHAQUE .tex =====
\documentclass[11pt,a4paper]{article}
\usepackage[utf8]{inputenc}\usepackage[T1]{fontenc}\usepackage{lmodern}
\usepackage[french]{babel}
\usepackage{amsmath,amssymb,mathtools}\usepackage{icomma}
\usepackage{siunitx}\sisetup{locale=FR}  % \qty{}{}, \si{}, \ang{}
\usepackage{xcolor}\usepackage{colortbl}
\usepackage{tikz}\usepackage{pgfplots}\pgfplotsset{compat=1.18}
\usepackage[siunitx,europeanresistors]{circuitikz} % circuits (élec.)
\usepackage[most]{tcolorbox}\usepackage{enumitem}\usepackage{array,booktabs}
\usepackage[a4paper,margin=2.2cm]{geometry}
\usetikzlibrary{arrows.meta,calc,angles,quotes,positioning,patterns,%
  backgrounds,shapes.geometric,decorations.pathmorphing,decorations.markings,intersections}
\definecolor{primary}{RGB}{222,49,99}\definecolor{primarydark}{RGB}{150,22,55}
\definecolor{defcol}{RGB}{0,114,178}\definecolor{methcol}{RGB}{52,92,148}
\definecolor{excol}{RGB}{0,135,90}\definecolor{attncol}{RGB}{200,40,55}
\tcbset{boxbase/.style={enhanced,breakable,boxrule=0.9pt,arc=2.5pt,
  left=3mm,right=3mm,top=2.3mm,bottom=2.3mm,fonttitle=\bfseries,coltitle=white}}
% --- boîtes TUTORIEL (noms EXACTS, ne pas renommer) ---
\newtcolorbox{cle}[1][]{boxbase,colback=defcol!6,colframe=defcol,
  title={\textsc{À retenir}\ifx\relax#1\relax\relax\else\textmd{ : \itshape #1}\fi}}
\newtcolorbox{methode}[1][]{boxbase,colback=methcol!7,colframe=methcol,sharp corners,
  title={\textsc{Méthode}\ifx\relax#1\relax\relax\else\textmd{ --- \itshape #1}\fi}}
\newtcolorbox{exemplebox}[1][]{boxbase,colback=excol!5,colframe=excol,
  title={\textsc{Exemple}\ifx\relax#1\relax\relax\else\textmd{ : \itshape #1}\fi}}
\newtcolorbox{piege}[1][]{boxbase,colback=attncol!6,colframe=attncol,
  title={\textsc{Piège}\ifx\relax#1\relax\relax\else\textmd{ : \itshape #1}\fi}}
\newtcolorbox{remarque}[1][]{boxbase,colback=black!4,colframe=black!45,
  title={\textsc{Remarque}\ifx\relax#1\relax\relax\else\textmd{ : \itshape #1}\fi}}
% --- boîtes EXERCICE / CORRIGÉ (noms EXACTS, auto-comptées) ---
\newtcolorbox[auto counter]{exo}[1][]{boxbase,colback=primary!4,colframe=primary,
  title={\textsc{Exercice}~\thetcbcounter\ifx\relax#1\relax\relax\else\textmd{ --- \itshape #1}\fi}}
\newtcolorbox[auto counter]{corrige}[1][]{boxbase,colback=excol!4,colframe=excol,
  title={\textsc{Corrigé}~\thetcbcounter\ifx\relax#1\relax\relax\else\textmd{ --- \itshape #1}\fi}}
\newcommand{\vv}[1]{\vec{#1}}\newcommand{\ds}{\displaystyle}
\newcommand{\R}{\mathbb{R}}\newcommand{\N}{\mathbb{N}}\newcommand{\Z}{\mathbb{Z}}
% (un \usepackage{fancyhdr}+en-tête est possible ; il est ignoré côté web)
% =========================================================================
\begin{document}
\sisetup{per-mode=symbol}

\begin{exo}[où le champ s'annule-t-il ?]
Deux charges positives $q_A = \SI{+1}{\micro\coulomb}$ et $q_B = \SI{+9}{\micro\coulomb}$ sont
distantes de $L = \SI{40}{\centi\metre}$.
\begin{center}
\begin{tikzpicture}[scale=1,>=Stealth,font=\footnotesize,
    ch/.style={circle,fill=attncol,minimum size=7mm,inner sep=0pt,text=white,font=\bfseries\scriptsize}]
  \draw[black!35] (-0.5,0)--(5.3,0);
  \node[ch] at (0,0) {$+$}; \node[below,font=\scriptsize] at (0,-0.45) {$q_A=1$};
  \node[ch,minimum size=9mm] at (4.8,0) {$+$}; \node[below,font=\scriptsize] at (4.8,-0.45) {$q_B=9$};
  \draw[<->,black!45,dashed] (0,0.6)--(4.8,0.6) node[midway,above,font=\scriptsize]{$L=\SI{40}{\centi\metre}$};
\end{tikzpicture}
\end{center}
\begin{enumerate}[label=\alph*)]
  \item Où le champ total peut-il s'annuler (entre, à gauche, à droite) ? Justifie.
  \item Calcule la distance de ce point à $q_A$.
\end{enumerate}
\end{exo}

\begin{exo}[le champ au milieu d'un dipôle]
Deux charges opposées $q_A = \SI{+1}{\micro\coulomb}$ (en $A$) et $q_B = \SI{-1}{\micro\coulomb}$
(en $B$) sont distantes de $\SI{20}{\centi\metre}$. On note $M$ le milieu de $[AB]$.
\begin{center}
\begin{tikzpicture}[scale=1,>=Stealth,font=\footnotesize,
    ch/.style={circle,minimum size=7mm,inner sep=0pt,text=white,font=\bfseries\scriptsize}]
  \draw[black!35] (-0.5,0)--(5.3,0);
  \node[ch,fill=attncol] at (0,0) {$+$}; \node[below,font=\scriptsize] at (0,-0.45) {$A$};
  \node[circle,fill=black,inner sep=1.2pt] at (2.4,0) {}; \node[below,font=\scriptsize] at (2.4,-0.4) {$M$};
  \node[ch,fill=defcol] at (4.8,0) {$-$}; \node[below,font=\scriptsize] at (4.8,-0.45) {$B$};
\end{tikzpicture}
\end{center}
\begin{enumerate}[label=\alph*)]
  \item Quel est le sens du champ $\vv{E}_A$ créé par $A$ en $M$ ? et celui de $\vv{E}_B$ ?
  \item Calcule l'intensité du champ total en $M$.
\end{enumerate}
\end{exo}

\begin{exo}[au centre d'un carré]
Quatre charges \emph{identiques} $+q$ sont placées aux quatre sommets d'un carré. On s'intéresse au
champ au centre $O$.
\begin{center}
\begin{tikzpicture}[scale=0.8,>=Stealth,font=\footnotesize,
    ch/.style={circle,fill=attncol,minimum size=7mm,inner sep=0pt,text=white,font=\bfseries\scriptsize}]
  \node[ch] (TL) at (-1.6,1.6) {$+q$}; \node[ch] (TR) at (1.6,1.6) {$+q$};
  \node[ch] (BL) at (-1.6,-1.6) {$+q$}; \node[ch] (BR) at (1.6,-1.6) {$+q$};
  \draw[black!30] (TL.center)--(TR.center)--(BR.center)--(BL.center)--cycle;
  \node[circle,fill=black,inner sep=1.2pt] at (0,0) {}; \node[font=\scriptsize,below right] at (0,0) {$O$};
\end{tikzpicture}
\end{center}
\begin{enumerate}[label=\alph*)]
  \item Que vaut le champ total en $O$ ? Justifie par la symétrie.
  \item On \emph{retire} la charge du coin haut-gauche. Vers quel coin pointe désormais le champ en
        $O$ ?
\end{enumerate}
\end{exo}

\begin{exo}[du champ à la force]
Une charge source $Q = \SI{+5}{\micro\coulomb}$ est placée en $O$. On étudie le point $M$ situé à
$r = \SI{30}{\centi\metre}$ de $O$.
\begin{enumerate}[label=\alph*)]
  \item Calcule l'intensité et le sens du champ $\vv{E}$ créé par $Q$ en $M$.
  \item On place en $M$ une charge $q = \SI{-2}{\nano\coulomb}$. Calcule l'intensité et le sens de
        la force qu'elle subit.
\end{enumerate}
\end{exo}

\begin{exo}[le champ ignore la charge test]
Au point $M$, une charge test $q = \SI{+2}{\nano\coulomb}$ subit une force $F = \SI{6e-4}{\newton}$.
\begin{enumerate}[label=\alph*)]
  \item Calcule le champ électrique $E$ en $M$.
  \item On remplace la charge test par $q' = 3q$. Quelle force subit-elle, et quel champ en
        déduit-on ? Que constate-t-on ?
\end{enumerate}
\end{exo}

\end{document}
